\section{Prolate coordinates and a prime-power application}
\label{v6:sec:coordinates}
\begin{proposition}[Full-band coordinates]\label{prop:T5a:involution}
Let $C_{2\tau}\psi_n=\ell_n\psi_n$ be the real orthonormal prolate
basis, with $\ell_n=\lambda_n(2\pi\tau)>0$ and parity $(-1)^n$.
Then $\chi_n=\ell_n^{-1/2}R_\tau U\psi_n=i^n\varphi_n$ form an
orthonormal basis of $L^2(-\tau,\tau)$, with $\varphi_n$ real of
parity $(-1)^n$. The matrix
$\mathcal J_{mn}=\langle\chi_m,J_\tau\chi_n\rangle$ is a real
symmetric involution, vanishes for opposite parities, and obeys
\begin{equation}\label{eq:T5a:matrix}
(\mathcal K_\tau)_{mn}=\sqrt{\ell_m\ell_n}\mathcal J_{mn},\quad
(C_\pm)_{mn}=\sqrt{\ell_m\ell_n}\,
                    (\delta_{mn}\pm\mathcal J_{mn})/2 .
\end{equation}
For equal parities,
$\mathcal J_{mn}=2(-1)^{(n-m)/2}
\int_0^\tau\varphi_m(\xi)\varphi_n(\xi-\tau)\,d\xi$.
All series in
\begin{equation}\label{eq:T5a:traces}
\begin{aligned}
\tr\mathcal K_\tau&=\sum_n\ell_n\mathcal J_{nn},&
\tr\mathcal K_\tau^2&=\sum_{m,n}\ell_m\ell_n\mathcal J_{mn}^2,\\
D_J&=\sum_{m,n}\ell_n(1-\ell_m)\mathcal J_{mn}^2,&
D_R&=\sum_n\ell_n(1-\ell_n)
\end{aligned}
\end{equation}
converge absolutely.
\end{proposition}
\begin{proof}
For $A=R_\tau U$, $A^*A=C_{2\tau}$ and both $A,A^*$ are
injective by the entire-function argument in
Theorem~\ref{thm:rev:structure:inertia}.
Thus the $\chi_n$ are orthonormal and complete: orthogonality of
$g$ to all $\chi_n$ implies $A^*g=0$.
Separating the sine and cosine parts of $U\psi_n$ gives $i^n\varphi_n$.
The exchange is a self-adjoint involution commuting with reflection;
its matrix entries are
$i^{n-m}\langle\varphi_m,J_\tau\varphi_n\rangle$, real for equal
parities and zero otherwise. Splitting the exchange integral into
two halves gives the stated formula.
Now $\mathcal K_\tau=A^*J_\tau A$ and
$C_\pm=A^*(I\pm J_\tau)A/2$ give \eqref{eq:T5a:matrix}.
Parseval yields $\sum_m\mathcal J_{mn}^2=1$; together with
$\sum_n\ell_n=4\tau$, this proves the trace formulas and all
absolute-convergence assertions.
\end{proof}

\subsection{One prime-power summand}\label{sec:arith}
For a prime power $q$, put $a_q=\Lam(q)/\sqrt q$, where $\Lam$
is the von Mangoldt function, and
$\varepsilon_q(c)=1-\log q/\log c$, $c>1$.
With the localized Weil-form basis $L^{-1/2}e^{2\pi inx/L}$,
$L=\log c$, $|n|\le N$, and prime distribution
$-\sum_{q\le c}a_q\delta_{\log q}$, the individual summand is
\begin{equation}\label{eq:5.1}
 -a_qA_N(\varepsilon_q(c)),\qquad
 \tau=N\varepsilon_q(c),\qquad
 \log(c/q)=\frac{\tau\log q}{N-\tau}.
\end{equation}
The matrix formula is \cite[Proposition~4.1]{CvS}, specialized as in
\cite{SilvaPrime}; the last equality is exact inversion.
\begin{corollary}[Finite-section design rule]\label{cor:design}
For fixed $q,T,\eta>0$, take $N>2T$, $N\ge(2T+8\pi T^2)/\eta$
and $c(\tau)=q^{N/(N-\tau)}$. Uniformly on $0<\tau\le T$, the embedded
normalized block is within $\eta$ of $\mathcal K_\tau$, and the physical
summand within $a_q\eta$ of $-a_q\mathcal K_\tau$.
If $\eta<\min(p_\tau,m_\tau)$ the normalized block has a positive even
and negative odd eigenvalue; multiplication by $-a_q$ reverses them.
\end{corollary}
\begin{proof}
The inversion gives $\varepsilon_q=\tau/N<1/2$; apply
Theorem~\ref{thm:nystrom} and \eqref{eq:3.7}.
\end{proof}
The cutoff layer has width $T\log q/(N-T)=O_{q,T}(N^{-1})$.
For phase-aligned eigenvectors at a fixed simple isolated parity eigenvalue,
Theorem~\ref{thm:packet} gives weights $1/2,1/2$ at the exchanged-arc
endpoints: for every $\zeta>0$, a sufficiently large fixed $C_\zeta$
makes each $C_\zeta/N$ neighborhood's limiting mass at least $1/2-\zeta$.
The disjoint-arc range is exactly $\sqrt c<q<c$, or $0<\tau<N/2$.
At $q=\sqrt c$ the matrix is $\operatorname{diag}((-1)^m)$; entrywise,
\begin{equation}\label{eq:5.6}
 A_N(1-\varepsilon)=-A_N(\varepsilon)
 +2\operatorname{diag}(\cos(2\pi m\varepsilon))_{m=-N}^N.
\end{equation}
Also $0<a_q\le2/e$, so signed count estimates transfer by reversing signs
and replacing a threshold $0<\alpha<a_q$ by $\alpha/a_q$.
This describes one moving summand, with no sign conclusion for the complete
Weil form. Its other prime and archimedean terms also move; fixed-mode
continuity alone gives no dimension-uniform right neighborhood.
The complete-form estimates \cite{Zhu2026,Angelone2026} and activation-wall
analysis \cite{Tao2026} concern different questions.
